\documentclass[11pt]{article}

\usepackage[margin=1in]{geometry}
\usepackage{amsmath}
\usepackage{amssymb}
\usepackage{booktabs}
\usepackage{enumitem}

\title{Mixed-Distance Tile Packing on a Checkerboard Chip:\\
Choosing a Code Distance per Site}
\author{qSNOW square-packing experiment}
\date{}

\newcommand{\Pset}{\mathcal{P}}
\newcommand{\Dset}{\mathcal{D}}
\newcommand{\Fp}{F}

\begin{document}
\maketitle

This note extends the single-distance model of \texttt{square\_packing\_model.tex}
to the case where each site may host a tile of \emph{any} distance drawn from a set
$\Dset$ (or none). The set is arbitrary and discovered from the data---in the
current experiments $\Dset = \{3, 5, 7\}$, with further distances to come---so the
model, the code, and this document are written for a general $\Dset$ and specialize
to any concrete choice. We reuse the base document's terminology without change:
\textbf{site} $(r,c)$ on the checkerboard sublattice ($r \equiv c \pmod 2$, unit
coordinate spacing), \textbf{physical qubit} (one per site), \textbf{tile}
(a distance-$d$ rotated-surface-code patch), and \textbf{origin} (the
minimum-coordinate corner of a tile's footprint). Only the geometry constants,
the conflict rule, the decision variables, and the constraints change to admit
several tile sizes; the reader is assumed familiar with Sections~1--3 of the base
document.

\section{Motivation}

A larger-distance tile generally has a lower logical error rate (LER) than a
smaller-distance tile at the same site, but occupies a larger footprint. Whether a
larger, more reliable tile is worth its area depends on the noise regime and on
the validity threshold $\tau$. In the low-noise data, tightening $\tau$ invalidates
small-distance placements faster than large-distance ones, so at strict $\tau$ only
the larger tiles remain valid over most of the chip; at looser $\tau$ the smaller
tiles win by packing more densely. The mixed model chooses, per site, which
distance (if any) to place so as to maximize the number of valid, non-overlapping
tiles.

\section{Geometry with several footprint sizes}
\label{sec:geometry}

\paragraph{Per-tile span.}
As in the base document, a distance-$d$ tile has footprint span $s_d = 2d + 1$
(so $s_3 = 7$, $s_5 = 11$, $s_7 = 15$), a geometric distance in coordinate units.
A tile of distance $d$ placed at origin $p = (r_p, c_p)$ occupies the footprint
\[
  \Fp(p, d) \;=\; \bigl\{\, (r, c) \;:\; r_p \le r \le r_p + s_d,\;\;
  c_p \le c \le c_p + s_d,\;\; r \equiv c \!\!\pmod 2 \,\bigr\},
\]
a square block whose side depends on $d$.

\paragraph{Asymmetric conflict rule.}
Consider two placements $a = (p, d_a)$ and $b = (q, d_b)$ with spans $s_a \equiv
s_{d_a}$ and $s_b \equiv s_{d_b}$. Their row windows $[r_p, r_p + s_a]$ and
$[r_q, r_q + s_b]$ intersect iff $r_p \le r_q + s_b$ and $r_q \le r_p + s_a$;
writing $\Delta r = r_p - r_q$ this is $-s_a \le \Delta r \le s_b$. Applying the
same reasoning to columns, and requiring overlap on \emph{both} axes, gives the
conflict relation
\begin{equation}
  a \sim b
  \quad\Longleftrightarrow\quad
  -s_a \le r_p - r_q \le s_b
  \;\;\text{and}\;\;
  -s_a \le c_p - c_q \le s_b .
  \label{eq:conflict}
\end{equation}
The bounds are \emph{asymmetric}: the upper bound uses $b$'s span and the lower
bound $a$'s. The relation is nonetheless symmetric in $a \leftrightarrow b$
(swapping negates $\Delta r$ and exchanges the bounds). When $s_a = s_b = s$ it
collapses to the base document's $|\Delta r| \le s$ and $|\Delta c| \le s$, so
\eqref{eq:conflict} is a strict generalization that applies to any pair of
distances in $\Dset$. Note two placements at the same origin
($\Delta r = \Delta c = 0$) always conflict, since $-s_a \le 0 \le s_b$: at most
one tile may sit at a given site.

\paragraph{Worked example.}
A $d=5$ tile at $(0,0)$ ($s_a = 11$) and a $d=3$ tile at $(0,8)$ ($s_b = 7$):
$\Delta c = -8$, and $-11 \le -8 \le 7$ holds, so they \emph{conflict} (windows
$[0,11]$ and $[8,15]$ overlap in columns). Reversing sizes---$d=3$ at $(0,0)$
($s_a=7$) and $d=5$ at $(0,8)$ ($s_b=11$)---gives $\Delta c = -8$ and the test
$-7 \le -8$ \emph{fails}, so they do \emph{not} conflict (windows $[0,7]$ and
$[8,19]$ are disjoint). The asymmetry is real and must not be replaced by
$|\Delta c| \le s$ for a single $s$.

\section{Candidates and sets}

Rather than a set of origins, the mixed model works with \emph{candidates}: a
candidate is a pair $(p, d)$ of an origin and a distance $d \in \Dset$ whose
measured LER $\ell_{p,d}$ satisfies the validity threshold. With a (possibly
per-distance) threshold $\tau_d$,
\[
  \Pset \;=\; \bigl\{\, (p, d) \;:\; d \in \Dset,\;\;
  \ell_{p,d} \le \tau_d \,\bigr\}.
\]
A single scalar $\tau$ is the common case ($\tau_d = \tau$ for all $d$); a stricter
$\tau_d$ for larger $d$ can be used to demand more of the larger tiles. Each chip
site thus offers a choice among the distances valid there.

\paragraph{Data note: nesting and monotonicity.}
Two structural facts about the experiment data simplify both the model and its
visualization; neither is assumed by the model, but both hold in practice.
\begin{itemize}[itemsep=2pt]
  \item \textbf{Origin nesting.} A larger distance's origin set is contained in a
        smaller distance's: the larger window cannot fit in the chip's edge band
        (the last $s_d - s_{d'} $ coordinate lines on each far side for $d > d'$),
        so edge sites offer only the smaller distances while the interior offers
        the full range. Concretely
        $\text{origins}(d_1) \supseteq \text{origins}(d_2) \supseteq \cdots$ for
        $d_1 < d_2 < \cdots$.
  \item \textbf{LER monotonicity.} Where both fit, a larger distance has the lower
        LER: $\ell_{p,d}$ is nonincreasing in $d$ at a fixed site $p$. Hence if
        $(p, d)$ is valid then so is $(p, d')$ for every larger $d' \in \Dset$ that
        fits at $p$. The valid distances at a site therefore form an \emph{upper
        set}, and the site is fully characterized by its \emph{smallest valid
        distance}---the cheapest (densest-packing) code that meets $\tau$ there.
        The visualization uses exactly this: each origin is drawn with the marker
        of its smallest valid distance, or a ``no valid tile'' marker if none fits.
\end{itemize}
Neither fact requires special handling---both follow automatically from reading
each distance's results file---and the model below remains correct even if a future
data set violates monotonicity, since it enumerates candidates explicitly rather
than assuming the upper-set structure.

\section{Decision variables}

For each candidate we introduce one binary variable:
\[
  y_{p,d} \in \{0,1\}, \qquad (p,d) \in \Pset,
  \qquad
  y_{p,d} =
  \begin{cases}
    1 & \text{if a distance-$d$ tile is placed at origin } p,\\
    0 & \text{otherwise.}
  \end{cases}
\]

\section{Model}

\begin{align}
  \max_{y} \quad & \sum_{(p,d) \in \Pset} y_{p,d}
    && \text{(maximize the number of placed tiles)} \label{eq:obj}\\[4pt]
  \text{s.t.} \quad
    & \sum_{d \,:\, (p,d) \in \Pset} y_{p,d} \le 1
    && \forall\, \text{origins } p
    && \text{(at most one tile per site)} \label{eq:origin}\\[2pt]
    & \sum_{(p,d) \in \mathcal{C}_u} y_{p,d} \le 1
    && \forall\, \text{sites } u
    && \text{(no two placed tiles overlap)} \label{eq:clique}\\[2pt]
    & y_{p,d} \in \{0,1\}
    && \forall\, (p,d) \in \Pset. && \label{eq:bin}
\end{align}

Objective~\eqref{eq:obj} counts placed tiles. Constraint~\eqref{eq:origin} forbids
placing more than one tile (of whatever distance) at the same site. Constraint
\eqref{eq:clique} is the non-overlap requirement in clique (site-cover) form,
defined next.

\paragraph{Generalized clique constraint.}
For a site $u = (u_r, u_c)$, let
\[
  \mathcal{C}_u \;=\; \bigl\{\, (p,d) \in \Pset \;:\;
  r_p \le u_r \le r_p + s_d \;\text{and}\; c_p \le u_c \le c_p + s_d \,\bigr\}
\]
be the candidates whose footprint covers $u$; note the membership test uses each
candidate's own span $s_d$. Any two candidates in $\mathcal{C}_u$ share $u$ and
hence conflict, so at most one may be used, giving \eqref{eq:clique}.

\paragraph{Validity and enumeration.}
Footprints remain axis-aligned boxes---only their sizes differ---so Helly's
theorem in the plane still applies: any set of pairwise-conflicting tiles has a
common covered site. Concretely, for any conflicting pair the point
$\bigl(\max(r_p, r_q),\, \max(c_p, c_q)\bigr)$ lies in both footprints. Both
coordinates are origin coordinates of candidates, so it suffices to enumerate the
witness site $u$ over the cross product of distinct origin rows and distinct
origin columns; each distinct membership set $\mathcal{C}_u$ with $|\mathcal{C}_u|
\ge 2$ yields a constraint, and keeping only the maximal sets suffices to cover
every conflicting pair. Each clique constraint dominates the pairwise edges among
its members, so \eqref{eq:clique} gives a tighter LP relaxation than the pairwise
form
$y_a + y_b \le 1 \;\; \forall\, a \sim b$, while describing the same integer set.
Constraint~\eqref{eq:origin} is in fact implied by \eqref{eq:clique} (co-located
candidates both cover the origin site and hence share a clique); it is retained
for clarity and as a guard. Nothing in the enumeration depends on $|\Dset|$: the
witness corner and the per-candidate membership test handle any collection of
distances uniformly.

\section{Behavior}
\begin{itemize}[itemsep=2pt]
  \item \textbf{Strict $\tau$ (low-noise regime).} Almost no small-distance
        placement is valid while many large-distance placements are; the optimum is
        dominated by the largest distances, whose footprint is the price of having
        any valid tile at all. (On the $30\times30$, $\texttt{mean\_0\_001}$ data
        at $\tau = 4\times10^{-4}$ the optimum places six $d=7$ and five $d=5$
        tiles and no $d=3$ tile.)
  \item \textbf{Loose $\tau$.} Every distance is broadly valid; the smallest tile
        packs most densely, so a count-maximizing optimum favors the smallest
        distance in $\Dset$.
  \item \textbf{Intermediate $\tau$.} A genuine mix arises---larger tiles in the
        quieter interior where they are valid, smaller tiles filling the edge band
        and the gaps between larger footprints.
\end{itemize}

\section{Notes}
\begin{itemize}[itemsep=2pt]
  \item The threshold-count objective~\eqref{eq:obj} is the direct generalization
        of the single-distance model. Other objectives fit the same variables and
        constraints---e.g.\ a reliability-weighted packing
        $\max \sum w_{p,d}\, y_{p,d}$, or placing at least $K$ tiles while
        minimizing total or worst-case LER (a quantity-versus-quality
        trade-off). Only the objective~\eqref{eq:obj} changes.
  \item Per-distance thresholds $\tau_d$ enter only through the candidate set
        $\Pset$; the constraints are unchanged.
  \item The distance set $\Dset$ enters only through $\Pset$ as well: adding a new
        distance means adding its results file and its candidates, with no change to
        the geometry, constraints, or enumeration.
  \item As in the base model, chip-boundary feasibility is baked into $\Pset$: the
        results files list only origins whose full footprint fits on the chip,
        which is exactly why the larger-distance edge bands are absent.
\end{itemize}

\end{document}
